\scanpage{32}
\diagram{\begin{tikzcd}
e \arrow[r,no head] & G_m^{\text{tordu}} \arrow[r] & G \arrow[r] & \mathbb Z/2 \arrow[r] \arrow[d,equal] & e \\
e \arrow[r] & \mathbb Z/3 \arrow[u] \arrow[r] & \mathfrak S_3 \arrow[r] & \mathbb Z/2 \arrow[r] & e
\end{tikzcd}}

On suppose tjrs qu'il n'y a pas de pt $s$ de car. 2 ou 3 dans $S$. Alors \deleted{$\unsure{\mathfrak S_3}$ opère} $\mathfrak S_3\times \unsure{Z|\mathbb Z}_2$ opère ``tamely'' sur $P^1_S$ [en fait, il est compatible avec la fibration $P^1_S\to S$, et opère ``tamely'' sur les fibres\ldots] Le quotient \inserted{par un groupe ou un ss-groupe} \deleted{\unsure{il est|qui est}} \inserted{\unsure{est isomorphe}} encore à $P^1_S$ [car il est \uline{lisse} sur $S$, et rationnel sur $S$ d'après Lüroth]. Les schémas de pts fixes $(0,1,\infty)$ $(2,-1,\tfrac12)$ $(\bar S=x',x'')$ --- qui sont disjoints, et chacun stable sous $\mathfrak S_3$, qui \deleted{\unsure{y}} \inserted{lui-même} opère de façon \deleted{\unsure{simplement}} transitive sur chaque fibre, donnent dans $P^1_S/\mathfrak S_3$ trois \uline{sections} distinctes [à expliciter\ldots] qu'on peut utiliser pour identifier $P^1_S/\mathfrak S_3$ à $P^1_S$ avec les sections $\infty,0,1$. De même, \deleted{\unsure{le schéma du revt}} $P^1_S/(\mathbb Z/3)$ est muni d'abord de deux sections qui correspondent à $(0,1,\infty)$ et $(2,-1,\tfrac12)$, et qui
\mspage{31}
